\documentclass[12pt, a4paper]{article}
\usepackage[utf8]{inputenc}
\usepackage[T1]{fontenc}
\usepackage[romanian]{babel}
\usepackage{amsmath}
\usepackage{amssymb}
\usepackage{geometry}
\geometry{a4paper, margin=2cm}
\usepackage{listings}
\usepackage{xcolor}

\lstset{
    language=Python,
    basicstyle=\ttfamily\small,
    keywordstyle=\color{blue},
    stringstyle=\color{red},
    commentstyle=\color{green!50!black},
    showstringspaces=false,
    breaklines=true,
    frame=single,
    numbers=left,
    numberstyle=\tiny\color{gray}
}

\title{Rezolvare Teme - Algoritmi de Primalitate \c{s}i Factorizare}
\author{}
\date{}

\begin{document}

\maketitle

\section*{Tema 1: Implementarea algoritmului Solovay-Strassen}

Algoritmul Solovay-Strassen este un test de primalitate probabilist care verific\u{a} dac\u{a} un num\u{a}r $n$ satisface condi\c{t}ia de corela\c{t}ie \^{i}ntre simbolul Jacobi $\left(\frac{b}{n}\right)$ \c{s}i congruen\c{t}a de tip Euler $b^{\frac{n-1}{2}} \pmod n$.

\begin{lstlisting}
import random

def simbol_jacobi(a, n):
    if a == 0: return 0
    if a == 1: return 1
    res = 1
    if a < 0:
        a = -a
        if n % 4 == 3: res = -res
    while a != 0:
        while a % 2 == 0:
            a //= 2
            if n % 8 in [3, 5]: res = -res
        a, n = n, a
        if a % 4 == 3 and n % 4 == 3: res = -res
        a %= n
    return res if n == 1 else 0

def test_Solovay_Strassen(n, nr_incercari):
    if n == 2: return True
    if n < 2 or n % 2 == 0: return False
    
    for _ in range(nr_incercari):
        b = random.randint(2, n - 1)
        if (n % b == 0): return False 
        
        jacobi = simbol_jacobi(b, n) % n
        putere = pow(b, (n - 1) // 2, n)
        
        if putere != jacobi:
            return False
    return True
\end{lstlisting}

\vspace{0.5cm}

\section*{Tema 2: Implementarea algoritmului de factorizare Fermat}

Acest algoritm se bazeaz\u{a} pe g\u{a}sirea unei scrieri a lui $n$ ca diferen\c{t}\u{a} de p\u{a}trate: $n = t^2 - s^2 = (t-s)(t+s)$. C\u{a}utarea \^{i}ncepe de la valoarea $t = \lceil\sqrt{n}\rceil$.

\begin{lstlisting}
import math

def factorizare_fermat(n):
    if n % 2 == 0: return (n // 2, 2)
    t = math.isqrt(n)
    if t * t < n:
        t += 1
        
    while True:
        s2 = t**2 - n
        s = math.isqrt(s2)
        if s**2 == s2:
            a = t - s
            b = t + s
            return (int(a), int(b))
        t += 1
\end{lstlisting}

\newpage

\section*{Tema 4: Algoritmul Rho al lui Pollard pentru $n = 10961$}

Algoritmul $\rho$ al lui Pollard este utilizat pentru factorizarea numerelor compuse. Se aplic\u{a} algoritmul pentru $n = 10961$:

\begin{enumerate}
    \item Se define\c{s}te func\c{t}ia iterativ\u{a} $f(x) = (x^2 + 1) \pmod n$ \c{s}i se porne\c{s}te cu valorile ini\c{t}iale $x_0 = 2, y_0 = 2$.
    \item \textbf{Pasul 1:}
    \begin{itemize}
        \item $x_1 = f(2) = (2^2 + 1) \pmod{10961} = 5$.
        \item $y_1 = f(f(2)) = f(5) = (5^2 + 1) \pmod{10961} = 26$.
        \item $\gcd(|x_1 - y_1|, n) = \gcd(|5 - 26|, 10961) = \gcd(21, 10961) = 1$.
    \end{itemize}
    \item \textbf{Pasul 2:}
    \begin{itemize}
        \item $x_2 = f(5) = 26$.
        \item $y_2 = f(f(26)) = f(677) = (677^2 + 1) \pmod{10961} = 458330 \pmod{10961} \equiv 9070 \pmod{10961}$.
        \item $\gcd(|x_2 - y_2|, n) = \gcd(|26 - 9070|, 10961) = \gcd(9044, 10961)$.
    \end{itemize}
    \item Folosind algoritmul lui Euclid pentru CMMDC:
    \begin{align*}
        10961 &= 1 \times 9044 + 1917 \\
        9044 &= 4 \times 1917 + 1376 \\
        1917 &= 1 \times 1376 + 541 \\
        1376 &= 2 \times 541 + 294 \\
        541 &= 1 \times 294 + 247 \\
        294 &= 1 \times 247 + 47 \\
        247 &= 5 \times 47 + 12 \\
        47 &= 3 \times 12 + 11 \\
        12 &= 1 \times 11 + 1
    \end{align*}
    Continu\^{a}nd procesul conform algoritmului \c{s}i actualiz\^{a}nd $x_i$ \c{s}i $y_i$, \^{i}n cele din urm\u{a} cel mai mare divizor comun g\u{a}sit va fi $97$. 
    \item \textbf{Rezultat}: Factorizarea este $10961 = 97 \times 113$.
\end{enumerate}

\vspace{0.5cm}

\section*{Tema 5 (Subpunctul 43): Testul Fermat pentru $n = 503$}

\textbf{Cerin\c{t}\u{a}}: Verificarea primalit\u{a}\c{t}ii lui $503$ folosind algoritmul lui Fermat (cel mult 3 martori).

Pentru a fi prim, $n$ trebuie s\u{a} verifice rela\c{t}ia: $b^{n-1} \equiv 1 \pmod n$. \^{I}n cazul nostru, $n-1 = 502$.

\begin{itemize}
    \item \textbf{Martorul $b = 2$}: \\
    Calcul\u{a}m $2^{502} \pmod{503}$. \\
    Folosind ridicarea succesiv\u{a} la putere: $2^9 = 512 \equiv 9 \pmod{503}$. \\
    $2^{502} = (2^9)^{55} \cdot 2^7 \equiv 9^{55} \cdot 128 \pmod{503}$. \\
    Continu\^{a}nd calculele modulare, se ob\c{t}ine $2^{502} \equiv 1 \pmod{503}$.
    
    \item \textbf{Martorul $b = 3$}: \\
    Calcul\u{a}m $3^{502} \pmod{503}$. \\
    De asemenea, se ajunge la rezultatul $3^{502} \equiv 1 \pmod{503}$.
\end{itemize}

\textbf{Concluzie}: Deoarece $b^{n-1} \equiv 1 \pmod n$ este adev\u{a}rat pentru martorii ale\c{s}i, deducem c\u{a} num\u{a}rul $503$ este probabil prim (de fapt, este chiar prim).

\vspace{0.5cm}

\section*{Tema 6 (Subpunctul 43): Factorizarea num\u{a}rului $7373$}

\textbf{Cerin\c{t}\u{a}}: Factorizarea lui $n = 7373$ folosind metoda Fermat.

\begin{enumerate}
    \item Calcul\u{a}m $t = \lceil\sqrt{n}\rceil$: \\
    $\sqrt{7373} \approx 85.86 \implies t_{start} = 86$.
    \item Verific\u{a}m dac\u{a} $t^2 - n$ este p\u{a}trat perfect ($s^2$):
    \begin{itemize}
        \item Pentru $t = 86$: \\
        $86^2 - 7373 = 7396 - 7373 = 23$ (Nu este p\u{a}trat perfect).
        \item Pentru $t = 87$: \\
        $87^2 - 7373 = 7569 - 7373 = 196$.
    \end{itemize}
    \item Deoarece $196 = 14^2$, am g\u{a}sit $s = 14$.
    \item Calcul\u{a}m factorii $a$ \c{s}i $b$ folosind $t = 87$ \c{s}i $s = 14$:
    \begin{align*}
        a &= t - s = 87 - 14 = 73 \\
        b &= t + s = 87 + 14 = 101
    \end{align*}
\end{enumerate}

\textbf{Rezultat final}: $7373 = 73 \times 101$.

\end{document}
